\subsection{Fibonacci--Lie Resonance Ladder: Interpreting Gauge Group Structures as Discrete Resonances of Resolution Windows}\label{subsec:fib-lie-resonance}

\paragraph{``Non-extraneous'' discrete anchor points for stable type dimensions}
Under the golden mean grammar (forbidding adjacent \(11\)), the stable type set of length \(m\) satisfies
\(
|X_m|=F_{m+2}
\).
Therefore, when one regards ``the number of generators/gauge bosons of gauge degrees of freedom'' as a discrete dimension interface that can be read out at a given layer,
a natural and auditable candidate is: at several small windows, \(|X_m|\) exactly equals the dimension of a compact Lie algebra,
thereby interpreting \(\mathrm{SU}(2)\), \(\mathrm{SU}(3)\), etc.\ as \emph{resonance points} in the resolution growth process, rather than extraneous assumptions.

\paragraph{Rigorous conclusion: Fibonacci--Lie resonance ladder (reproducible exhaustive search)}
Table~\ref{tab:fibonacci_lie_resonance_ladder} lists all exact resonance points found in the exhaustive sweep for \(m\le 500\):
\input{sections/generated/tab_fibonacci_lie_resonance_ladder}
Moreover, more strongly, if the equality is written directly as a quadratic Diophantine equation, then within the same exhaustive range only very few solutions appear:
\input{sections/generated/eq_fib_lie_resonance_solutions}
Both outputs above are generated by the script \texttt{scripts/exp\_fibonacci\_lie\_resonance\_ladder.py}
and constitute auditable, reproducible arithmetic facts (independent of any interpretive-layer narrative).

\begin{corollary}[Scarcity of Fibonacci--Lie resonance points and the arithmetic minimality of \texorpdfstring{$\mathfrak{su}(2),\mathfrak{su}(3)$}{su2,su3}]\label{cor:fib-lie-resonance-scarcity-su2-su3}
Within the exhaustive sweep range \(m\le 500\), the resonance points satisfying
\(
|X_m|=F_{m+2}=\dim(\mathfrak g)
\)
where \(\mathfrak g\) is a compact simple Lie algebra (or its equidimensional dual type) occur only at
\[
m\in\{2,4,6,8\},
\qquad
F_{m+2}\in\{3,8,21,55\}
\]
(Table~\ref{tab:fibonacci_lie_resonance_ladder}).
Among these, \(3\) and \(8\) uniquely correspond in the compact simple classification to \(\mathfrak{su}(2)\) and \(\mathfrak{su}(3)\), respectively;
therefore, under this discrete arithmetic interface, \(\mathfrak{su}(2)\oplus\mathfrak{su}(3)\) as the non-abelian visible part manifests as the two earliest nontrivial resonances.
\end{corollary}
\begin{proof}
The list of resonance points follows directly from Table~\ref{tab:fibonacci_lie_resonance_ladder}; the uniqueness of the compact simple isomorphism classes with \(\dim=3,8\) is a standard classification fact. \qed
\end{proof}

\begin{remark}[Four-step shift with the boundary word tower]\label{rem:fib-lie-resonance-shift-by-4}
In Subsection~\ref{subsec:bdry-tower-zeck-gut}, we use the cardinality of the endpoint \((1,1)\) fiber
\(
|X_m^{\mathrm{bdry}}|=F_{m-2}
\)
as the ``boundary mode count.''
Note that
\[
|X_m^{\mathrm{bdry}}|=F_{m-2}=F_{(m-4)+2}=|X_{m-4}|,
\]
so the resonance ladder of this subsection (based on \(|X_m|\)) manifests in the boundary tower caliber as a discrete shift \(m\mapsto m+4\):
the stable type resonance points \(m=2,4,6,8,\dots\) correspond to boundary layers \(m=6,8,10,12,\dots\).
This explains why the boundary three-layer tower (\(m=4,6,8\)) naturally aligns with the SM's \(1,3,8\) three-segment gauge boson counting interface:
it is a readout of the same Fibonacci dimension skeleton on different fiber cross-sections.
\end{remark}

\begin{remark}[Minimal extra generator count: \(21-12=9\)]\label{rem:fib-lie-extra-9}
If the stable type dimension \( |X_6|=21 \) at \(m=6\) is regarded as the candidate generator count for the ``minimal unification layer,''
and the visible boson count of the SM is taken as \(12\) (the boundary tower interface of Subsection~\ref{subsec:bdry-tower-zeck-gut}),
then the gap between the two is a rigid integer:
\[
21-12=9.
\]
This provides a counting fingerprint distinct from the ``choose a large group first, then decompose'' a priori group decomposition approach: extra channels must appear at the discrete scale of \(9\) generators,
and their specific realization should be audited through this framework's twisted spectrum/boundary layer discrete labels/projection gate mechanisms (rather than introduced as continuously tunable parameters).
\end{remark}

\begin{remark}[The \texorpdfstring{$(\ZZ_2)^{21}$}{(Z2)\^21} binary charge lattice of the fold-gauge (\texorpdfstring{$m=6$}{m=6}, dyadic readout)]\label{rem:fold-gauge-z2-21}
Beyond the ``dimension layer'' with stable type dimension \(|X_6|=21\), the \(m=6\) window also yields an independent, auditable charge lattice at the fiber redundancy layer.
For the binary interval folding \(\mathrm{Fold}^{\mathrm{bin}}_6:\{0,\dots,63\}\to X_6\) (Definition~\ref{def:fold-bin}), its fiber automorphism group
\[
\mathcal{G}^{\mathrm{bin}}_6:=\mathrm{Aut}\!\bigl(\mathrm{Fold}^{\mathrm{bin}}_6\bigr)
\]
satisfies
\[
\bigl(\mathcal{G}^{\mathrm{bin}}_6\bigr)^{\mathrm{ab}}\cong (\ZZ_2)^{21}
\]
(Corollary~\ref{cor:fold-bin-gauge-m6}).
Here each \(\ZZ_2\) factor can be read out via the sign representation of the permutation within the corresponding fiber, interpretable as ``the minimal nontrivial binary gauge charge visible at the smallest nontrivial window.''
This \(21\)-dimensional binary charge lattice matches in dimension the \(B_3/C_3\) seed dictionary of \S\ref{subsec:bdry-tower-zeck-gut}, and provides an additional evidence chain for promoting geometric interpretations such as
\(
21=\dim\wedge^2\RR^7=\dim\mathfrak{so}(7)
\)
to auditable constraints.
\end{remark}
